% !TeX root = ../main.tex
\section{Context-free languages}

\subsection{Context-free grammars, formally}

A \textbf{context-free grammar} (CFG) is exactly the grammar definition from file 02,
restricted so every left-hand side is a single nonterminal:

\begin{definition}
A grammar $G = (N, \Sigma, P, S)$ is \textbf{context-free} if every production in $P$ has
the form $A \to \beta$, with $A \in N$ a single nonterminal and $\beta \in (N \cup
\Sigma)^*$ arbitrary (including $\varepsilon$).
\end{definition}

``Context-free'' names the restriction precisely: you're allowed to replace $A$ regardless
of what surrounds it in the sentential form --- the replacement doesn't depend on
\emph{context}. This single restriction is what makes CFGs efficiently parseable
(polynomial time, in general; linear time for the well-behaved subclasses --- LL, LR --- a
real parser generator targets), while still being expressive enough for essentially every
programming language's syntax. Everything from file 02 (derivations, leftmost/rightmost,
$L(G)$, ambiguity) applies unchanged; this file is about the specific tools and pathologies
that come with the CFG restriction.

\subsection{Parse trees and derivation order}

A \textbf{parse tree} makes a derivation's structure explicit as a tree: the root is $S$,
each internal node is a nonterminal with children given by the right-hand side of the rule
used to expand it, and the leaves, read left to right, spell out the derived string.

The key fact worth internalizing: \emph{leftmost and rightmost derivations of the same
string, using the same grammar, produce the same parse tree} --- they just build it in a
different order (leftmost expands the tree top-down, left child first; rightmost expands
right child first). This is why parser implementations bother distinguishing them:
recursive-descent (top-down) parsers naturally produce a leftmost derivation, while
shift-reduce (bottom-up) parsers naturally produce a rightmost derivation \emph{in reverse}.
Different algorithms, same tree, same language --- that equivalence is exactly what lets you
choose a parsing strategy based on implementation convenience rather than which trees you
want to allow.

\subsection{Ambiguity in CFGs, revisited: fixing dangling-else}

File 02 showed the dangling-else grammar is ambiguous:
\[
S \to \mathbf{if}\ b\ \mathbf{then}\ S \mid \mathbf{if}\ b\ \mathbf{then}\ S\ \mathbf{else}\ S \mid \mathbf{other}
\]
Now let's actually fix it, rather than just naming the problem. The standard convention
(``else binds to the nearest unmatched if'') can be built directly into the grammar by
distinguishing \emph{matched} statements (every \texttt{if} inside has its \texttt{else})
from \emph{unmatched} ones (some \texttt{if} is still waiting for an \texttt{else}):

\[
\begin{aligned}
S &\to M \mid U \\
M &\to \mathbf{if}\ b\ \mathbf{then}\ M\ \mathbf{else}\ M \mid \mathbf{other} \\
U &\to \mathbf{if}\ b\ \mathbf{then}\ S \mid \mathbf{if}\ b\ \mathbf{then}\ M\ \mathbf{else}\ U
\end{aligned}
\]

The rule that does all the work is the last one: an unmatched if-then-else must have a
\emph{matched} statement between \texttt{then} and \texttt{else} (you can't leave an
unmatched if dangling in there --- any \texttt{if} before the \texttt{else} must already
have found its own \texttt{else}), and only the \emph{tail} $U$ after the \texttt{else} is
allowed to be unmatched. This forces every \texttt{else} to attach to the innermost
\texttt{if} that doesn't already have one.

\begin{example}
Take $\mathbf{if}\ b_1\ \mathbf{then}\ \mathbf{if}\ b_2\ \mathbf{then}\ S_1\ \mathbf{else}\
S_2$ (using $S_1, S_2$ as stand-ins for \texttt{other}). Under the new grammar this string
is derivable only as an unmatched statement $U$, and only in one way:
\[
U \to \mathbf{if}\ b_1\ \mathbf{then}\ S
  \to \mathbf{if}\ b_1\ \mathbf{then}\ U
  \to \mathbf{if}\ b_1\ \mathbf{then}\ (\mathbf{if}\ b_2\ \mathbf{then}\ M\ \mathbf{else}\ U)
\]
with $M \to \mathbf{other}$ ($S_1$) and $U \to \mathbf{other}$ ($S_2$). There is no
alternative derivation that attaches the \texttt{else} to the outer \texttt{if}, because
doing so would require the outer $\mathbf{if}\ b_1\ \mathbf{then}\ M\ \mathbf{else}\ U$
rule, which demands the \emph{then}-branch be a fully matched $M$ --- but
$\mathbf{if}\ b_2\ \mathbf{then}\ S_1$ (with no \texttt{else}) is not a matched statement.
The grammar itself rules out the second parse tree. That's the difference between ``we
adopt the convention that else binds to the nearest if'' (a rule bolted onto an ambiguous
grammar) and ``the grammar is unambiguous and only generates the intended structure''
(the actual fix).
\end{example}

\subsection{Chomsky Normal Form}

Some algorithms that operate on CFGs (notably the CYK parsing algorithm) need the grammar
in a very restricted shape first.

\begin{definition}
A CFG is in \textbf{Chomsky Normal Form} (CNF) if every production has one of two forms:
$A \to BC$ (two nonterminals) or $A \to a$ (a single terminal) --- plus, optionally,
$S \to \varepsilon$ if $\varepsilon \in L(G)$, with $S$ not appearing on the right of any
rule.
\end{definition}

\begin{theorem}
Every context-free language is generated by some grammar in Chomsky Normal Form.
\end{theorem}

The conversion algorithm runs in a few mechanical passes:
\begin{enumerate}
  \item \textbf{Eliminate $\varepsilon$-productions} ($A \to \varepsilon$) other than
    possibly $S \to \varepsilon$, by inlining: for every rule containing $A$ on the right,
    add a version with $A$ removed (handling the case where $A$ appears multiple times by
    considering all subsets to remove).
  \item \textbf{Eliminate unit productions} ($A \to B$, a single nonterminal), by
    replacing $A \to B$ with a copy of every rule $B \to \gamma$ as $A \to \gamma$.
  \item \textbf{Break long right-hand sides}: any $A \to X_1 X_2 \cdots X_k$ with $k > 2$
    gets split using fresh nonterminals: $A \to X_1 A_1$, $A_1 \to X_2 A_2$, \ldots,
    $A_{k-2} \to X_{k-1} X_k$.
  \item \textbf{Replace terminals mixed with nonterminals}: any rule with a terminal $a$
    alongside another symbol (e.g.\ $A \to aB$) introduces a fresh nonterminal
    $T_a \to a$ and rewrites as $A \to T_a B$.
\end{enumerate}

Why bother: CNF guarantees every derivation of a string of length $n$ takes exactly
$2n - 1$ steps (each step either splits one nonterminal into two, or terminates one into a
single terminal), which is exactly the uniform structure the CYK dynamic-programming parser
needs to fill an $O(n^2)$ table of ``which nonterminals can derive this substring'' in
$O(n^3)$ time overall. Without that uniformity, there's no clean way to bound the table-filling
recursion.

\subsection{Pushdown automata}

A DFA/NFA (file 03) has no memory beyond its current state, which is precisely why it can't
recognize $a^nb^n$ (file 04). A \textbf{pushdown automaton} (PDA) fixes this by adding a
single unbounded stack --- the minimal amount of extra memory that gets you from regular to
context-free.

Informally: a PDA reads input left to right, like an NFA, but at each step it may also pop
a symbol off the stack, and push zero or more symbols onto it, and its transitions can be
conditioned on what's currently on top of the stack. Acceptance is either by ending in an
accepting state, or (an equivalent convention) by emptying the stack.

\begin{definition}
A \textbf{pushdown automaton} is a 7-tuple $M = (Q, \Sigma, \Gamma, \delta, q_0, Z_0, F)$
where $Q, \Sigma, q_0, F$ are as for an NFA, $\Gamma$ is the \textbf{stack alphabet},
$Z_0 \in \Gamma$ is the initial stack symbol, and
$\delta: Q \times (\Sigma \cup \{\varepsilon\}) \times \Gamma \to \mathcal{P}(Q \times
\Gamma^*)$ maps a state, an input symbol (or $\varepsilon$), and the top-of-stack symbol to
a set of (next state, string to push in place of the popped symbol) pairs.
\end{definition}

\begin{example}
A PDA for $\{a^nb^n \mid n \ge 0\}$: push an $A$ marker for every $a$ read, then pop one
marker for every $b$ read, and accept exactly when the stack is back to just $Z_0$ and
input is exhausted.
\[
\delta(q_0, a, Z_0) = \{(q_0, AZ_0)\}, \qquad
\delta(q_0, a, A) = \{(q_0, AA)\}
\]
\[
\delta(q_0, b, A) = \{(q_1, \varepsilon)\}, \qquad
\delta(q_1, b, A) = \{(q_1, \varepsilon)\}
\]
with $q_0$ the start state ($q_0 \in F$, to accept $n=0$) and $q_1 \in F$ accepting once all
the $A$'s pushed for the $a$'s have been popped by matching $b$'s. Reading $aabb$: push
$A$, push $A$ (stack: $AAZ_0$, top to bottom), then pop $A$ on the first $b$ (state
switches to $q_1$), pop $A$ on the second $b$; stack is back to $Z_0$, input exhausted,
accept. Try $aaab$: after three pushes the stack has three $A$'s, but only one $b$ pops one
$A$, leaving two $A$'s on the stack when input runs out --- not accepted, correctly.
\end{example}

\begin{theorem}
A language is context-free if and only if some pushdown automaton accepts it.
\end{theorem}

We state this without proof (the constructions in both directions --- CFG-to-PDA via a
generic ``simulate a leftmost derivation using the stack to hold the pending sentential
form'' recipe, and PDA-to-CFG via a construction with nonterminals indexed by pairs of
states --- are standard but mechanical). The practical upshot mirrors file 03's NFA/DFA
story exactly: PDAs are the ``machine'' view of context-free languages, CFGs are the
``generative'' view, and they describe the same class.

\subsection{The pumping lemma for context-free languages}

Just as DFAs have bounded memory (finitely many states), PDAs have a stack whose
\emph{depth} is bounded by a function of the parse tree's height for CNF grammars --- which
gives a pumping lemma with \emph{two} pumped pieces instead of one.

\begin{theorem}[Pumping lemma for CFLs]
If $L$ is context-free, there is a constant $p \ge 1$ such that every $z \in L$ with
$|z| \ge p$ can be split $z = uvwxy$ with:
\begin{enumerate}
  \item $|vx| \ge 1$ (at least one of $v, x$ is nonempty),
  \item $|vwx| \le p$,
  \item $uv^iwx^iy \in L$ for every $i \ge 0$.
\end{enumerate}
\end{theorem}

The proof idea: take a CNF grammar for $L$ with $p'$ nonterminals; any parse tree of height
$> p'$ must, by pigeonhole, repeat some nonterminal $A$ along its longest root-to-leaf path.
The subtree between the two occurrences of $A$ produces the ``middle'' pumped piece $v \dots
x$, and it can be spliced in or cut out repeatedly because both occurrences of $A$ derive
in the same way.

\subsubsection{Using it: proving $\{a^nb^nc^n \mid n \ge 0\}$ is not context-free}

Same adversary structure as file 04's regular pumping lemma, just with two pieces to track
instead of one.

\begin{example}
Let $L = \{a^nb^nc^n \mid n \ge 0\}$. Suppose $L$ is context-free with pumping length $p$.

\textbf{Choose the string.} Pick $z = a^pb^pc^p \in L$, $|z| = 3p \ge p$.

\textbf{The split is out of your control}, but condition 2 constrains it usefully: since
$|vwx| \le p$, the substring $vwx$ cannot span all three blocks ($a$'s, $b$'s, $c$'s) at
once --- $p$ symbols aren't enough to reach from the first $a$-block into the $c$-block,
since the $b$-block alone has length $p$ and sits between them. So $vwx$ touches \emph{at
most two} of the three distinct letters.

\textbf{Pump it.} Take $i = 2$: $uv^2wx^2y$ increases the total length of the letters
touched by $v$ and $x$ combined, while condition 1 guarantees $|vx| \ge 1$ so \emph{some}
letter's count strictly increases. Because $vwx$ touches at most two of $\{a,b,c\}$, at
least one letter's count is left completely unchanged by pumping, while at least one other
letter's count strictly grows. So $uv^2wx^2y$ no longer has equal numbers of $a$'s, $b$'s,
and $c$'s --- contradicting condition 3.

This holds for every possible split satisfying conditions 1--2, so no valid $p$ exists, and
$L$ is not context-free. $\blacksquare$
\end{example}

Intuitively: a PDA's single stack can enforce \emph{one} equality constraint (match $a$'s
against $b$'s, using the stack as a counter) but has no second independent counter left over
to also match against $c$'s. This is exactly the sense in which context-free languages are
one step more powerful than regular languages, but still bounded.

\subsection{Closure properties of CFLs}

\begin{proposition}
Context-free languages are closed under union, concatenation, and Kleene star.
\end{proposition}

All three constructions are direct grammar surgery: given CFGs $G_1, G_2$ (with disjoint
nonterminal sets, and start symbols $S_1, S_2$) for $L_1, L_2$, add a fresh start symbol $S$
and:
\begin{itemize}
  \item \textbf{Union}: $S \to S_1 \mid S_2$.
  \item \textbf{Concatenation}: $S \to S_1 S_2$.
  \item \textbf{Star}: $S \to S_1 S \mid \varepsilon$.
\end{itemize}
Each is a one-line addition to the grammar, which is why CFLs are so easy to combine
compared to what's coming next.

\begin{proposition}
Context-free languages are \textbf{not}, in general, closed under intersection or
complement.
\end{proposition}

\begin{example}
Let $L_1 = \{a^nb^nc^m \mid n, m \ge 0\}$ and $L_2 = \{a^mb^nc^n \mid n, m \ge 0\}$. Both
are context-free: $L_1$'s grammar matches $a$'s against $b$'s with a stack and then just
appends arbitrarily many $c$'s ($S \to AC$, $A \to aAb \mid \varepsilon$, $C \to cC \mid
\varepsilon$); $L_2$ is symmetric, matching $b$'s against $c$'s and prefixing arbitrary
$a$'s. But
\[
L_1 \cap L_2 = \{ a^nb^nc^n \mid n \ge 0 \},
\]
exactly the language we just proved is \emph{not} context-free. So the intersection of two
CFLs need not be a CFL --- closure under intersection fails, and this single example is the
standard counterexample cited for exactly this fact.
\end{example}

Non-closure under intersection then immediately forces non-closure under complement too, by
De Morgan: if CFLs were closed under complement, they'd also be closed under intersection
via $L_1 \cap L_2 = \overline{\overline{L_1} \cup \overline{L_2}}$ (union is closed, so this
identity would make intersection closed too) --- contradicting the example above. This is
the mirror image of file 03's regular-language closure properties: regular languages got
all three (union, intersection, complement) essentially for free via the product
construction; CFLs only get the ``easy'' three (union, concatenation, star) and lose the
other two, because a single stack can't be duplicated and run ``in parallel'' against a
second stack the way two DFAs' finite states can.

\subsection{Forward pointer: this is exactly what parsers handle}

Context-free grammars are, almost without exception, the grammar class used to describe
real programming language syntax --- arithmetic expressions, statements, function
definitions, all of it. The next block of this course is about turning a CFG into a working
parser: recursive-descent (top-down, LL) and shift-reduce (bottom-up, LR) algorithms that
build the parse tree efficiently, the precedence and ambiguity-removal techniques from this
file and file 02 applied for real, and what happens when a language's true syntax rules
(``this identifier must already be declared'') fall outside what any CFG can express at
all --- pushed, as mentioned back in file 02, into a semantic-analysis phase that runs after
the parser has already built its tree.

\subsection{Exercises}

\begin{itemize}
  \item Convert the grammar $S \to aSb \mid ab \mid \varepsilon$ into Chomsky Normal Form,
    showing each of the four passes (eliminate $\varepsilon$-productions, eliminate unit
    productions, break long right-hand sides, isolate terminals).
  \item Draw the two distinct parse trees for
    $\mathbf{if}\ b_1\ \mathbf{then}\ \mathbf{if}\ b_2\ \mathbf{then}\ S_1\ \mathbf{else}\
    S_2$ under the \emph{original ambiguous} dangling-else grammar from file 02, then
    confirm by hand that only one of those two trees is derivable under the fixed
    $M$/$U$ grammar in this file.
  \item Extend the $a^nb^n$ PDA in this section to a PDA for $\{a^nb^{2n} \mid n \ge 0\}$
    (hint: push two markers per $a$, or otherwise adjust how many pops each $b$ performs).
  \item Use the CFL pumping lemma to show $\{ a^ib^jc^k \mid i < j < k \}$ is not
    context-free. Follow the adversary structure explicitly.
  \item $L_1 \cap L_2$ can fail to be context-free even when $L_1, L_2$ are, but what if
    one of the two languages is \emph{regular} instead of merely context-free? Look up (or
    conjecture and then check) whether CFLs are closed under intersection with a regular
    language, and sketch why the product-construction idea from file 03 does or doesn't
    carry over.
  \item Is every context-free language also context-sensitive (Type 1 in the Chomsky
    hierarchy from file 02)? Justify your answer by relating the CNF production shapes to
    the context-sensitive production shape $\alpha A \beta \to \alpha\gamma\beta$.
\end{itemize}
